% !TeX program = xelatex
% FORMAT_ID: academic-presentation-uva
% FORMAT_VARIANT: core-template
% Reading edition with a physical-page count and progress line.
% Dark edition with the programme beneath UvA and institutional marks on the page ground.
\documentclass[aspectratio=169,11pt]{beamer}
% Local editors can resolve the private package from a project dependency directory.
\InputIfFileExists{.texinputs/ac-editor-inputs.tex}{}{}
\usepackage[dark,illc-artwork]{ac-illc-presentation}[2026/11/26]

\ACLogoRightsBasis{ILLC wordmark, CC0, illc-mol-thesis-template (ae51780e).}

\title[Forcing with Urelements]{Forcing with\\Urelements}
\subtitle{Witness names and Yao's repair}
\author{Simone Testino}
\ACauthorurl{https://github.com/SimoneTestino}
\ACprogramme{Master of Logic}{https://msclogic.illc.uva.nl/programme/mol/}
\ACsupervisor{Yurii Khomskii}{https://www.illc.uva.nl/People/person/3457/Dr-Yurii-Khomskii}
\ACsemester{Second semester 2025--2026}
\date[]{Last updated: 26 November 2026}
\ACsource{Bokai Yao}{Yao}{Set Theory with Urelements}
  {https://arxiv.org/pdf/2303.14274v3}{arXiv v3 (2023)}
\ACmaterials{%
  \ACgithublink{https://github.com/SimoneTestino/Academia}{GitHub}\quad
  \ACsourcelink{https://github.com/SimoneTestino/Academia/blob/main/Forcing_with_Urelements/Presentation/Forcing_with_Urelements.tex}\par
  \ACprojectlink{https://simonetestino.pages.dev/research/forcing-with-urelements/first-step/\#learning-cells}{Checked Isabelle/HOL exercise}}
\ACappendixname{Appendix}
\hypersetup{pdftitle={Forcing with Urelements: Witness names and Yao's repair},
  pdfauthor={Simone Testino},pdfsubject={Based on Bokai Yao, Set Theory with Urelements, arXiv v3}}

\begin{document}
\ACtitleframe

\ACsection{A forcing refresher}{}
\begin{frame}[t]{Atoms, sets and the empty set}
The language contains membership $\in$ and a primitive atom predicate $\mathcal A$.
\[
\mathcal A(a)\ \Longrightarrow\ \forall x\;(x\notin a),
\qquad \operatorname{Set}(x)\ \equiv\ \neg\mathcal A(x).
\]
\begin{block}{Having no members does not determine an object}
Two distinct atoms $a,b$ and the set $\varnothing$ can all have no members.
Extensionality identifies \emph{sets} with the same members.
\end{block}
Yao's base theory requires neither the existence of an atom nor a set of all atoms.
\medskip\par
Our example adds the explicit assumption $\mathcal A(a)$, $\mathcal A(b)$ and $a\ne b$.
\ACpages{\S1.2, Axiom A and Definition 1}{4--5}{11--12}{11}
\end{frame}

\begin{frame}[t]{The ground model and the forcing order}
\begin{ACcolumns}
\ACdiagramcolumn
\centering\ACanchoredbranches{1}{p}{q}{terminal}{terminal}
\ACtextcolumn
The \emph{ground model} $M$ satisfies set theory with urelements ($\mathrm{ZFU}_{R}$). In $M$ take distinct atoms $a,b$ and a pure, three-condition poset $\mathbb P=\{1,p,q\}$. Its only strict comparisons are $p<1$ and $q<1$.
\end{ACcolumns}
\begin{block}{Two assumptions about $M$}
\emph{Transitive} means that if $x\in M$ and $y\in x$, then $y\in M$.\par
\emph{Countable} means that, from outside $M$, we can enumerate its objects. The enumeration need not belong to $M$.
\end{block}
\ACpages{\S\S3.1--3.2 and Remark 55, specialised to three conditions}{44--48}{51--55}{51}
\end{frame}

\begin{frame}[t]{Which conditions can a filter contain?}
\begin{ACcolumns}
\ACdiagramcolumn
\centering\ACanchoredbranches{1}{p}{q}{left choice}{right choice}
\ACtextcolumn
\begin{block}{Filter}
A filter is a nonempty $G\subseteq\mathbb P$ with both properties.
\par\smallskip
\emph{Upward closure.} If $r\in G$ and $r\le s$ in $\mathbb P$, then $s\in G$.
\par\smallskip
\emph{Directedness.} For all $r,s\in G$ there is a $t\in G$ with $t\le r$ and $t\le s$.
\end{block}
\end{ACcolumns}
\smallskip\par
Upward closure puts the weaker condition $1$ in every filter. Since $p\perp q$, no condition strengthens both, so $\{1,p,q\}$ fails directedness. The filters are $\{1\}$, $\{1,p\}$ and $\{1,q\}$.
\smallskip\par
The \emph{poset atoms} $p,q$ are minimal conditions. The witnesses $a,b$ are urelements.
\ACpages{\S3.2 and Remark 55, worked finite example}{45--48}{52--55}{52}
\end{frame}

\begin{frame}[t]{Why there are exactly two generic filters}
\begin{ACcolumns}
\ACdiagramcolumn
\centering\ACanchoredbranches{1}{p}{q}{$G_p=\{1,p\}$}{$G_q=\{1,q\}$}
\ACtextcolumn
\begin{block}{Density and genericity}
A subset $D\subseteq\mathbb P$ is \emph{dense} if
\par\smallskip
\centerline{$\forall r\in\mathbb P\;\exists s\in D\;(s\le r).$}
\par\smallskip
A filter $G$ is \emph{$M$-generic} if $G\cap D\ne\varnothing$ for every dense $D\in M$.
\end{block}
\end{ACcolumns}
\medskip\par
\begin{ACcolumns}
\ACdiagramcolumn
\textbf{Testing $D=\{p,q\}$}\par\smallskip
$\begin{array}{@{}l@{\quad}ccc}
\text{start at }r & 1 & p & q\\
\text{choose }s & p & p & q
\end{array}$\par\smallskip
Here $s$ depends on $r$, and equality is allowed.
\ACtextcolumn
Since $\{p,q\}\in M$ is dense, a generic must contain $p$ or $q$. Directedness rules out both at once.
\smallskip\par
Each terminal condition is its own only strengthening. Every dense set contains both, so $G_p$ and $G_q$ are generic.
\end{ACcolumns}
\ACpages{\S3.2 and Remark 55, worked finite example}{45--48}{52--55}{52}
\end{frame}

\begin{frame}[t]{The question on two branches}
\begin{ACcolumns}
\ACdiagramcolumn
\centering\ACanchoredbranches{1}{p}{q}{set $\{a\}$\\member $a$}{set $\{b\}$\\member $b$}
\ACtextcolumn
Suppose the set we get is $\{a\}$ in $G_p$ and $\{b\}$ in $G_q$. Once the branch is known, we can point to its only member.
\medskip\par
A forcing \emph{name} $\sigma$ is chosen in $M$. Its value after a filter $G$ is chosen is written $\sigma_G$.
\end{ACcolumns}
\begin{block}{Can we choose one name before we know the branch?}
We need the \emph{same} $\sigma$ to give $\sigma_{G_p}=a$ and $\sigma_{G_q}=b$.
\end{block}
Both filters already belong to $M$, so here $M[G]=M$. What varies is the value of a name inside that model.
\ACpages{Remark 55 and the old valuation, finite instance}{47--48}{54--55}{54}
\end{frame}

\ACsection{Why old names fail}{Can one old name pick the member on both branches?}
\begin{frame}[t]{How to read an old name}
In Yao's old system, each atom $a$ is a name with $a_G=a$ on every branch. A \emph{set-name} $\sigma$ is a set of ordered pairs.
\[
\Bigl\langle\underbrace{\tau}_{\text{an old name}},\;
                \underbrace{r}_{\text{a condition}}\Bigr\rangle
\]
\begin{block}{To evaluate $\sigma$ at $G$}
Keep the pairs with $r\in G$. Evaluate each first entry at $G$ and collect its value as a member of the resulting set.
\par\smallskip
\centerline{$\displaystyle \sigma_G=\bigl\{\tau_G \bigm| \exists r\in G\;(\langle\tau,r\rangle\in\sigma)\bigr\}.$}
\end{block}
For example, $\langle a,p\rangle$ contributes $a$ whenever $p\in G$. If it is the only pair kept, the outer braces give $\sigma_G=\{a\}$.
\ACpages{Definition 51 and the old valuation}{46--47}{53--54}{53}
\end{frame}

\begin{frame}[t]{Evaluating the same set-name twice}
Choose one old set-name $s_{\#}=\{\langle a,p\rangle,\langle b,q\rangle\}$.
\medskip\par
\begin{ACcomparison}{@{}lccc@{}}
Filter $G$ & Pair kept & Value of $s_{\#}$ & Its member \\
\midrule
$G_p=\{1,p\}$ & $\langle a,p\rangle$ & $\{a\}$ & $a$ \\
$G_q=\{1,q\}$ & $\langle b,q\rangle$ & $\{b\}$ & $b$ \\
\end{ACcomparison}
\medskip\par
On the left, $p\in G_p$ and $q\notin G_p$. The one value to collect is $a_{G_p}=a$, giving $(s_{\#})_{G_p}=\{a\}$. On the right, only $q$ survives the test, giving $\{b\}$.
\medskip\par
Both filters contain $1$, and both values have a member. In this example that is exactly what $1\Vdash_{\#}\exists y\,(y\in s_{\#})$ means.
\ACpages{old valuation, Theorem 53 and Remark 55}{47--48}{54--55}{54}
\end{frame}

\begin{frame}[t]{Local witnesses and a uniform name}
Write $M^{\mathbb P}_{\#}$ for the old names in $M$. On each branch there is a witness.
\[
\forall G\in\{G_p,G_q\}\;\exists y\;(y\in(s_{\#})_G).
\]
For example, the old name $a$ works on $G_p$ but fails on $G_q$, since $a\notin\{b\}$.
\begin{block}{The uniform question quantifies over names first}
\[
\exists\sigma\in M^{\mathbb P}_{\#}\;\forall G\in\{G_p,G_q\}\;
\bigl(\sigma_G\in(s_{\#})_G\bigr)?
\]
\end{block}
The object $y$ in the first formula can depend on $G$. In the second, one $\sigma$ must give $\sigma_{G_p}=a$ and $\sigma_{G_q}=b$.
\ACpages{Remark 55; quantifier comparison for the finite example}{47--48}{54--55}{54}
\end{frame}

\begin{frame}[t]{Why every old witness name fails}
Suppose an old name $\sigma$ works on both branches. There are only two cases.
\begin{enumerate}
\item If $\sigma$ is an atom, then $\sigma_{G_p}=\sigma_{G_q}=\sigma$.
It cannot equal both $a$ and $b$, since $a\ne b$.
\item If $\sigma$ is a set-name, both values are sets.
Neither value can equal the urelement required on its branch.
\end{enumerate}
\begin{block}{The failure is a failure of fullness}
\[
1\Vdash_{\#}\exists y\;(y\in s_{\#}),\qquad
\nexists\sigma\in M^{\mathbb P}_{\#}\;1\Vdash_{\#}\sigma\in s_{\#}.
\]
\end{block}
Here forcing at $1$ means truth on both generic branches. The old forcing theorem and truth lemma still hold.
\ACpages{Theorems 53--54 and Remark 55}{47--48}{54--55}{54}
\end{frame}

\ACsection{Yao's repair and mixing}{How can one name take two different atom values?}
\begin{frame}[t]{Two roles for pairs in repaired names}
Every repaired name is a \emph{set} of pairs $\langle y,r\rangle$.
Here $y$ is either a urelement or another repaired name.
\begin{block}{The first component determines the role}
\begin{tabular}{@{}ll@{}}
$y=a$ is an atom & proposes $a$ as the value of the whole name\\[7pt]
$y=\sigma$ is a name & proposes $\sigma_G$ as a member of its set-value
\end{tabular}
\end{block}
The condition $r$ still controls when the proposal is active.
\medskip\par
Atoms themselves are no longer names. This opens the possibility of a set-name whose value is an atom.
\ACpages{Definition 57}{48}{55}{55}
\end{frame}

\begin{frame}[t]{Compatibility of repaired names}
\begin{block}{Compatibility requirement}
For two pairs in one repaired name $\tau$,
\[
\left.
\begin{gathered}
\langle a,r\rangle,\langle y,s\rangle\in\tau,\\
\mathcal A(a),\quad a\ne y
\end{gathered}
\right\}\ \Longrightarrow\ r\perp s.
\]
\end{block}
\begin{itemize}
\item Different atom-values must be attached to incompatible conditions.
\item An atom-value and a proposed member must also be attached to incompatible conditions: a name $y$ is a set, so $y\ne a$.
\end{itemize}
A filter can activate at most one atom-value. When it does, it activates no member proposal.
\ACpages{Definition 57, clause (ii)}{48}{55}{55}
\end{frame}

\begin{frame}[t]{Evaluating a repaired name}
\begin{block}{Valuation at $G$}
\[
\tau_G=
\begin{cases}
a,&\text{if }\langle a,r\rangle\in\tau,\ \mathcal A(a),\ r\in G,\\[4pt]
\{\sigma_G:\langle\sigma,r\rangle\in\tau,\ \sigma\text{ a name},\ r\in G\},
&\text{if no atom proposal is active.}
\end{cases}
\]
\end{block}
The compatibility rule makes the active atom unique and rules out active member proposals in that case.
\medskip\par
If no atom proposal is active, valuation collects the active members in the familiar way. If no pair is active, the value is $\varnothing$.
\ACpages{Definition 63, using Definition 57}{48, 53}{55, 60}{60}
\end{frame}

\begin{frame}[t]{Names for an atom and its singleton}
\begin{ACcolumns}
\ACequalcolumn
\textbf{A name for the atom}
\[
\check a=\{\langle a,1\rangle\}.
\]
The first component is the atom $a$.
Its pair proposes the whole value.
\[
(\check a)_G=a.
\]
\ACequalcolumn
\textbf{A name for its singleton}
\[
\check{\{a\}}=\{\langle\check a,1\rangle\}.
\]
The first component is the name $\check a$.
Its pair proposes a member.
\[
(\check{\{a\}})_G=\{a\}.
\]
\end{ACcolumns}
\medskip
Replacing $a$ by $\check a$ changes the role of the pair. The braces around a name alone do not tell us whether its value is an atom or a set.
\ACpages{Definitions 57 and 63}{48, 53}{55, 60}{55}
\end{frame}

\begin{frame}[t]{A repaired witness on both branches}
\[
\tau=\{\langle a,p\rangle,\langle b,q\rangle\},\qquad
s=\{\langle\check a,p\rangle,\langle\check b,q\rangle\}.
\]
\begin{ACcolumns}
\ACdiagramcolumn
\centering\ACbranchdiagram{1}{p}{q}{$\tau_{G_p}=a$\\$s_{G_p}=\{a\}$}{$\tau_{G_q}=b$\\$s_{G_q}=\{b\}$}
\ACtextcolumn
$p\perp q$ makes $\tau$ a valid repaired name.
The name $s$ has only member proposals.
\medskip\par
On either branch, $\tau_G\in s_G$. Hence
\[
1\Vdash\tau\in s.
\]
The displayed set of pairs for $\tau$ is the same as for $s_{\#}$, but its semantics has changed.
\end{ACcolumns}
\ACpages{Definitions 57 and 63; finite instance of Lemma 60}{48--50, 53}{55--57, 60}{55}
\end{frame}

\begin{frame}[t]{Mixing extends this construction}
Let $U^{\mathbb P}$ be the repaired names. On the antichain $D=\{p,q\}$, set
\[
f(p)=\check a,\qquad f(q)=\check b.
\]
The name $\tau$ is a mixture: $p\Vdash\tau=\check a$ and $q\Vdash\tau=\check b$.
\begin{block}{Yao's Mixing Lemma, Lemma 60}
For a given set function $f:D\to U^{\mathbb P}$, with $D$ an antichain, there is a repaired name $\mu$ such that
\[
\forall r\in D\qquad r\Vdash\mu=f(r).
\]
\end{block}
Mixing starts with a set function $f$ and needs neither Collection nor Choice. Gathering witness names is a separate step.
\ACpages{Lemma 60, statement and construction}{49--50}{56--57}{56}
\end{frame}

\begin{frame}[t]{Structure, renaming and representation}
Swap $a$ and $b$ consistently, leaving our pure poset fixed.
The same calculations give the same existence and impossibility results, with the atom labels exchanged.
\begin{block}{A structuralist reading of the example}
The phenomenon depends on the relation between branches, names and distinct witnesses, rather than on the labels of the atoms.
\end{block}
Here both constructions give $M[G]=M$ and use the same atoms. Changing the naming language changes which branch-dependent values admit one uniform name.
\medskip\par
Yao extends permutations of a set of atoms to membership automorphisms. The interpretation in this slide is our reading of the example, not a philosophical theorem attributed to Yao.
\ACpages{Permutations of atoms; Remark 55; Definition 57}{6, 47--48}{13, 54--55}{13}
\end{frame}

\ACsection{Collection and fullness}{When does every forced existential have a witness name?}
\begin{frame}[t]{Atoms, classes and gathering witnesses}
The formula $\mathcal A(x)$ defines the class of all atoms. It need not define a set.
\begin{itemize}
\item If some set contains all atoms, Separation gives the set of all atoms. That set may be infinite.
\item If no set contains them all, the class of atoms is proper, relative to the model.
\end{itemize}
\begin{block}{The gathering question is more specific}
For each input $x$ in a set $w$, suppose some witness $y$ exists.
Can one set contain a suitable witness for every $x$?
\end{block}
A proper class of atoms does not by itself refute this principle. The two-branch example needs only two distinct atoms; the general gathering problem is additional.
\ACpages{Definition 1; Definition 3 and Proposition 5}{5--6}{12--13}{12}
\end{frame}

\begin{frame}[t]{Replacement and Collection}
For a set $w$ and a fixed parameter $u$, Yao writes:
\medskip\par
\textbf{Replacement}
\[
\bigl(\forall x\in w\;\exists!y\;\varphi(x,y,u)\bigr)
\ \Longrightarrow\
\bigl(\exists v\;\forall x\in w\;\exists y\in v\;\varphi(x,y,u)\bigr).
\]
\textbf{Collection}
\[
\bigl(\forall x\in w\;\exists y\;\varphi(x,y,u)\bigr)
\ \Longrightarrow\
\bigl(\exists v\;\forall x\in w\;\exists y\in v\;\varphi(x,y,u)\bigr).
\]
The same set $v$ serves every $x$. It need not be unique and need not contain all possible witnesses. The formula's remaining parameters stay fixed.
\medskip\par
Each displayed implication is universally quantified over $w,u$, and each admissible formula gives an instance of the scheme.
\ACpages{\S1.2: Replacement at the end of p.~4; Collection at the start of p.~5}{4--5}{11--12}{11}
\end{frame}

\begin{frame}[t]{A common container for witnesses}
Let $w=\{1,2\}$ and let $\varphi(x,y)$ mean $y=10x$ or $y=10x+1$.
\[
\begin{array}{c|c|c}
x&\text{possible witnesses}&\text{one witness in }v=\{10,21\}\\\hline
1&10,11&10\\[3pt]
2&20,21&21
\end{array}
\]
\begin{itemize}
\item Every $x$ has a witness, so Collection's antecedent holds.
\item Witnesses are not unique, so Replacement's antecedent fails for this formula.
\item The displayed container satisfies the common conclusion.
\end{itemize}
A false antecedent does not make the Replacement implication false. This example illustrates the schemes; it does not separate their strength.
\ACpages{\S1.2; numerical illustration of the two schemes}{4--5}{11--12}{11}
\end{frame}

\begin{frame}[t]{Which scheme implies which?}
Write $U$ for unique witnesses, $E$ for existing witnesses and $C$ for a common witness container.
\[
U\Longrightarrow E,\qquad
\text{Replacement: }U\Longrightarrow C,\qquad
\text{Collection: }E\Longrightarrow C.
\]
\begin{block}{Collection implies Replacement}
Assume $U$. Then $E$ holds. Collection gives $C$.
Uniqueness was assumed, not inferred from Collection.
\end{block}
In pure ZF, Replacement also proves Collection using bounds on witness ranks.
In Yao's setting, Replacement need not prove Collection, even with Choice. With a proper class of atoms, a rank bound alone need not provide a set container.
\ACpages{Definition 1 and following paragraph; Proposition 5}{5--6}{12--13}{12}
\end{frame}

\begin{frame}[t]{Fullness for a general forcing poset}
\begin{block}{Fullness}
For each condition $p$, each formula and each tuple of name parameters,
\[
p\Vdash\exists y\;\varphi(y,\vec\tau)
\quad\Longrightarrow\quad
\exists\sigma\in U^{\mathbb P}\;p\Vdash\varphi(\sigma,\vec\tau).
\]
\end{block}
The name $\sigma$ may depend on $p$ and the formula. The condition $p$ stays fixed.
\medskip\par
The existential forcing clause gives witness names densely below $p$. Fullness asks for one name working already at $p$.
\medskip\par
The repaired two-branch construction proves this one instance directly. It does not establish fullness for arbitrary formulas or posets.
\ACpages{Definition 58(8); Theorem 61}{49--52}{56--59}{56}
\end{frame}

\begin{frame}[t]{Yao's general equivalence}
The base theory explicitly includes Choice:
\[
\mathrm{ZFU}_{R}=\mathrm{ZU}+\text{Replacement},\qquad
\mathrm{ZFCU}_{R}=\mathrm{ZFU}_{R}+\mathrm{AC}.
\]
Here $\mathrm{ZU}$ is Zermelo theory with atoms and Foundation.
\begin{block}{Theorem 61, over $\mathrm{ZFCU}_{R}$}
\[
\text{Collection}\quad\Longleftrightarrow\quad
\begin{gathered}
\text{for every forcing poset }\mathbb P,\\
\text{the repaired forcing relation is full.}
\end{gathered}
\]
\end{block}
The quantifier over \emph{all} posets matters. This is a general theorem about the repaired machinery, separate from the finite example.
\ACpages{Definition 1; Theorem 61}{5, 50--52}{12, 57--59}{57}
\end{frame}

\begin{frame}[t]{Collection, Choice and mixing}
Suppose $p\Vdash\exists y\;\varphi(y)$, with fixed name parameters.
\begin{enumerate}
\item Choice gives a maximal antichain $D$ of conditions below $p$ admitting local witness names.
\item Collection puts enough of those names into one set.
\item Choice selects a name $f(r)$ for each $r\in D$.
\item Mixing combines the given function $f$ into one name $\mu$.
\item Density below $p$ gives $p\Vdash\varphi(\mu)$.
\end{enumerate}
\begin{block}{The roles are distinct}
Collection gathers names. Choice selects them. Mixing combines them.
\end{block}
For the finite example, the two local names were already explicitly available. The general argument must first obtain a set of names.
\ACpages{Theorem 61, forward implication}{50--51}{57--58}{57}
\end{frame}

\begin{frame}[t]{Fullness, truth and preservation}
\begin{description}
\item[Fullness] A forced existential has a witness name under the \emph{same} condition.
\item[Truth lemma] Truth in one extension corresponds to forcing by \emph{some} condition in its generic filter:
\[
M[G]\models\varphi(\vec\tau_G)
\ \Longleftrightarrow\
\exists r\in G\;M\models r\Vdash\varphi(\vec\tau).
\]
\item[Preservation] Axioms holding in the ground model continue to hold in the extension, under the relevant theorem's assumptions.
\end{description}
The old names have a truth lemma despite failing fullness. The repaired forcing theorem over $\mathrm{ZFU}_{R}$ does not assume Choice or Collection.
\ACpages{Theorems 53, 61 and 66; preservation in \S3.2.4}{47, 50--62}{54, 57--69}{61}
\end{frame}

\begin{frame}[t]{Conclusions and the next question}
The old names describe the singleton on either branch, but cannot name its varying atom-member uniformly. Yao's repaired names can do both.
\medskip\par
Mixing generalises the repair. Over $\mathrm{ZFCU}_{R}$, Collection gives exactly the additional gathering strength needed for fullness for all posets.
\medskip\par
\textbf{Source.} Bokai Yao, \href{https://arxiv.org/pdf/2303.14274v3}{\emph{Set Theory with Urelements}}, doctoral dissertation, University of Notre Dame, arXiv v3.
\medskip\par
\textbf{Further work.} The \href{https://isa-afp.org/entries/Forcing.html}{Archive of Formal Proofs} develops forcing over pure ZFC in Isabelle/ZF. Our proposed finite development would compare the two valuations in the presence of atoms.
\medskip\par
\ACgithublink{https://github.com/SimoneTestino/Academia/tree/main/Forcing_with_Urelements}{Project materials}\qquad
\ACprojectlink{https://github.com/SimoneTestino/Academia/tree/main/Forcing_with_Urelements/Formalisation}{Isabelle project scope}
\end{frame}

\ACappendix{Technical details and next steps}

\begin{frame}[t]{Appendix: the mixing construction}
For a given set function $f:D\to U^{\mathbb P}$ on an antichain, define
\[
\mu=\bigcup_{r\in D}
\{\langle y,t\rangle:t\in\mathbb P,\ \exists s\;(\langle y,s\rangle\in f(r)\ \land\ t\leq r,s)\}.
\]
Pairs from $f(r)$ are restricted to common strengthenings with their assigned branch $r$.
\begin{itemize}
\item Within one branch, the original name's compatibility condition prevents conflicts.
\item Across different branches, the antichain prevents conflicts.
\item Under a branch $r\in D$, the resulting name agrees with $f(r)$.
\end{itemize}
Thus $r\Vdash\mu=f(r)$ for every $r\in D$.
\ACpages{Lemma 60, proof}{49--50}{56--57}{56}
\end{frame}

\begin{frame}[t]{Appendix: ranks and available atoms}
For a \emph{set} $A$ of atoms, the hierarchy $V_\alpha(A)$ has set-sized levels.
Proposition 5 proves restricted Collection: if every input has a witness in $V(A)$, one common set contains sufficient witnesses.
\begin{block}{Why Replacement suffices in that restricted situation}
Assign each input the least ordinal level containing a witness.
Replacement gathers these ordinals; a sufficiently high level gives the container.
\end{block}
For pure ZF use $A=\varnothing$. For a proper class of atoms, there need not be one set $A$ accommodating all required witnesses.
\medskip\par
Here $\operatorname{ker}(x)$ is the set of atoms in the transitive closure of $\{x\}$. In the reverse direction of Theorem 61, the witness name's kernel together with the poset's kernel bounds the atoms needed. Proposition 5 gives Collection.
\ACpages{Definition 2; Proposition 5; Theorem 61, reverse implication}{5--6, 51--52}{12--13, 58--59}{13}
\end{frame}

\begin{frame}[t]{Appendix: what the extension preserves}
For a countable transitive $M\models\mathrm{ZFU}_{R}$ and an $M$-generic $G$ over set forcing $\mathbb P\in M$:
\[
\begin{array}{c|c}
\text{ground model}&\text{extension}\\\hline
\mathrm{ZFU}_{R}&\mathrm{ZFU}_{R}\\[3pt]
\mathrm{ZFCU}_{R}&\mathrm{ZFCU}_{R}\\[3pt]
\mathrm{ZFCU}_{R}+\text{Collection}&\mathrm{ZFCU}_{R}+\text{Collection}
\end{array}
\]
There are no new atoms or ordinals. New sets of existing atoms may appear.
The kernel bound $\operatorname{ker}(\tau_G)\subseteq\operatorname{ker}(\tau)$ limits which atoms a name can use.
\medskip\par
These are preservation results, not a reformulation of fullness.
\ACpages{Lemmas 64--65; Theorems 67, 71 and 73}{53, 56--60}{60, 63--67}{60}
\end{frame}

\begin{frame}[t]{Appendix: when forcing really adds a set}
For comparison, let $\mathbb P=2^{<\omega}$ be finite binary strings, ordered by extension.
\[
\langle0,1,0\rangle\leq\langle0,1\rangle\leq\langle0\rangle.
\]
For each $n$, strings of length at least $n$ form a dense set. A generic filter determines
\[
c=\bigcup G:\omega\longrightarrow\{0,1\}.
\]
For each ground-model binary sequence $r$, conditions disagreeing with $r$ somewhere form a dense set in $M$. Thus $c\ne r$ for every such $r$.
\medskip\par
This is an extension adding a new set. Our terminal two-branch poset was chosen to test the naming language without that additional complication.
\ACpages{Classical forcing background to \S3.2; comparison example}{45--46}{52--53}{52}
\end{frame}

\begin{frame}[t]{Current work and next steps}
A finite Isabelle/HOL teaching exercise now checks the old-name obstruction and the repaired selector on this three-condition order.
\begin{block}{What is checked in the finite model}
\begin{itemize}
\item \textbf{Old names.} No single old name takes value $a$ in $G_p$ and $b$ in $G_q$.
\item \textbf{Yao's repair.} The selector $w=\{\langle a,p\rangle,\langle b,q\rangle\}$ is admissible, with $w_{G_p}=a$ and $w_{G_q}=b$.
\end{itemize}
\end{block}
My study and independent research development are ongoing. The general Collection/fullness theorem remains mathematical reading.
\smallskip\par
\ACprojectlink{https://simonetestino.pages.dev/research/forcing-with-urelements/first-step/\#learning-cells}{Checked teaching exercise}\qquad
\ACprojectlink{https://github.com/SimoneTestino/Academia/blob/main/Forcing_with_Urelements/Formalisation/README.md}{Research-code status}
\smallskip\par
\ACsourcelink{https://github.com/SimoneTestino/Academia/blob/main/Forcing_with_Urelements/Presentation/Forcing_with_Urelements.tex}\qquad
\href{https://arxiv.org/pdf/2303.14274v3}{Bokai Yao's full thesis}
\par
\emph{Thank you for your attention.}
\vfill\par
{\scriptsize\color{ACMuted}Student project at the University of Amsterdam and ILLC. \href{https://simonetestino.pages.dev/contact/}{Contact the author.}}
\end{frame}

\end{document}
